%=============================================================================
\section{Discussion}
\label{sec:discussion}
%=============================================================================

\subsection{What This Paper Establishes}

Four things, the first three in decreasing order of confidence and the fourth
the one that answers the objection this work is most exposed to.

That \emph{classical information gain systematically solicits the questions
users answer worst, in catalogues of a particular shape}, is established by
direct measurement of selector behaviour and does not depend on any assumed
$\varepsilon$. On the music catalogue, a classical selector asks questions
whose annotated crossover averages $0.225$ against a catalogue average of
$0.098$, filling $53.9\%$ of a twenty-question budget with the tier a listener
judges least reliably, while a uniformly random selector faces $0.097$, which
is the catalogue average. Whatever one believes about the numeric values of
$\varepsilon$, the ordering of the annotation is what drives that contrast, and
the contrast is large.

That \emph{the effect is manufactured by discretisation rather than inherited
from the domain} is established by three independent lines that agree:
Proposition~\ref{prop:discretisation}, which shows quantile binning fixes
apparent informativeness at $\Hb(1/k)$ irrespective of semantics; the
measurement that binned attributes carry mean maximum information gain $0.861$
against $0.618$ for natively categorical ones; and the observation that
refining the binning to four equal-population bins \emph{strengthens} the
measured bias, from $\rho = -0.44$ to $\rho = -0.61$. The absence of the effect
in two catalogues with no continuous attributes to bin is the prediction this
account makes, and it holds.

That \emph{correcting the acquisition function helps, and does so from an
annotation cheap enough to be deployable}, is established at $n=900$ paired
trials with exact tests and multiplicity control, and---more importantly---is
shown to survive the assumption it rests on: over $60$ worlds drawn from priors
on the tier values, the fixed deployed system wins in $60$ of $60$, with a
minimum advantage of $+1.33\pp$. Three coarse tiers recover roughly $72\%$ of
what perfect per-attribute knowledge would deliver, and a two-tier annotation
recovers $78\%$ of the three-tier gain (Section~\ref{sec:results-ablation}).

That \emph{the three assumed constants are not load-bearing} is answered three
ways, in increasing strength: the conclusion survives
a distribution over those constants; it survives a majority of attributes being
placed in the wrong tier; and the constants are not required at all. EM over
$600$ logged sessions of the already-deployed incumbent---no annotation, no
labels, and the target item latent throughout---recovers a per-attribute
crossover vector that matches the curator's annotation to a dead heat, $38$
discordant pairs in each direction out of $900$
(Section~\ref{sec:results-em}).

Two of the model's structural assumptions were attacked directly, and the
outcomes point in opposite directions. Relaxing independence across turns
\emph{strengthens} the result: making subjective errors persist for the rest of
a session costs classical selection $0.78\pp$ and costs reliability-aware
selection exactly nothing, because a selector that avoids unreliable attributes
is not exposed to their errors being persistent
(Section~\ref{sec:results-correlated}). Adding abstention, by contrast, refutes
a prediction we made in advance: we expected the penalties to compound and the
absolute advantage to grow, and it shrinks
(Section~\ref{sec:results-erasure}). Only the relative advantage compounds. We
report the failed prediction because a paper that only tests the predictions it
expects to confirm is not testing anything.

\subsection{What It Does Not Establish}

We claim no algorithmic novelty. The acquisition function is the mutual
information between the target and the observed answer under a binary symmetric
channel; it is in Cover and Thomas~\cite[Ch.~7]{cover2006}, and a reader who
finds \eqref{eq:raig} unsurprising is reading it correctly. The contribution is
that the observation model in some deployed attribute catalogues is
heterogeneous in a specific and adversarial way, that the heterogeneity is
partly an artefact of preprocessing, that a learned policy trained without a
noise model inherits the blind spot, and that the correction needs far less
knowledge about users than one would expect. Those are empirical claims, and
they are what the experiments are for.

Nor have we measured $\varepsilon$ on people, and it is worth being precise
about what the log-based estimator of Section~\ref{sec:results-em} does and
does not settle. It establishes identifiability and sufficiency: a
per-attribute crossover vector can be recovered from six hundred sessions of
ordinary use without annotation, labels or ground truth, and recovering it is
enough to match a hand-annotated one exactly. It does not establish that real
users behave this way, because the logs it consumes were generated by a
simulator whose users err at rates drawn from the annotation---so what EM
recovers is, at bottom, the structure we put in. The same caveat applies to
the ensemble over tier priors (Section~\ref{sec:results-ensemble}), the
corruption sweep (Section~\ref{sec:results-corruption}) and the repeated-item
measurement (Section~\ref{sec:reliability-empirical}). Together they establish
that the conclusion is not an artefact of three particular numbers. None of
them tells a practitioner what $\varepsilon$ is in their domain, and only the
study of Appendix~\ref{app:protocol} would.

Finally, the two negative results on catalogue-derived specifications should
not be softened, because together they say something the positive results do
not. RAIG-empirical, which infers reliability from the rate at which repeated
items disagree with each other, is beaten by the three-tier annotation by
$6.33\pp$ and is no better than classical information gain. RAIG-psycho, which
computes the crossover of each binned attribute from the position of its own
cut points, fails at every setting of its two parameters
(Section~\ref{sec:results-psycho}). Both measure a property of the
\emph{catalogue}, and both are defeated by the same thing: a listener cannot
name the key of a song, not because the boundary between keys is fuzzy or
because the catalogue disagrees with itself about it, but because they never
learned to name keys. Answerability is a fact about people. What succeeds---the
curator's tiers, and the vector learned from logs---succeeds because each is,
in the end, an observation of people rather than of data.

\subsection{A Deployment Test}

The cross-domain results yield a procedure a practitioner can run before
committing to anything, and it is the most transferable output of this work.

\begin{enumerate}
\item \emph{Is there anything to correct?} Compute per-attribute maximum
one-versus-rest entropy and compare it across a coarse answerability
annotation. A Kruskal--Wallis test on three tiers is the right instrument; a
Spearman correlation on three tied levels is not. If informativeness
concentrates on the hard-to-answer tier, classical selection is being actively
misled.
\item \emph{Is the correction affordable in information?} Compute the
identifiability ceiling of the reliable attribute subset by counting distinct
attribute vectors. This costs one database query. If the ceiling is high
(dermatology, $70.8\%$), reliability-aware selection will help substantially
even when step~1 finds no bias. If it is near zero (mushroom, $0.3\%$), the
unreliable attributes are load-bearing and no selection rule can route around
them.
\item \emph{Is the channel actually noisy?} Reliability weighting is a
correction for a noisy channel and is harmful when the channel is clean: in our
clean condition RAIG-tiered loses to classical selection by more than
fifty points at $T=20$, because it breaks ties towards attributes that cannot
by themselves separate the catalogue. Switch it off, not down, when answers are
reliable.
\item \emph{Annotate coarsely and ship, then let the system re-estimate.} Two
tiers recover $78\%$ of what three tiers deliver, and three tiers recover about
$72\%$ of what per-attribute truth would deliver. Once the system is running,
that residual is addressable without a measurement campaign: EM over its own
logs recovers a per-attribute vector that matches the annotation
(Section~\ref{sec:results-em}). The order of operations is therefore curator
first, logs second, user study only if the first two disagree.
\end{enumerate}

Steps~1 and~2 are computable from the catalogue alone. Step~3 requires knowing
roughly how noisy users are, which is the one thing that genuinely needs
observation---and Section~\ref{sec:results-lambda} locates the threshold.

\subsection{Threats to Validity}

\emph{Simulated users.} Every session in this paper is played by a simulator
drawing from a binary symmetric channel. This is the central threat and we do
not minimise it. Its severity is bounded in three ways: the deployed system
never sees the true $\varepsilon$; the conclusion is reported across a
distribution of true $\varepsilon$ rather than at a point; and two structural
assumptions of the channel---that every turn is an independent draw, and that a
question is always answered---are attacked directly in
Sections~\ref{sec:results-correlated} and~\ref{sec:results-erasure}. Three
things remain unaddressed. Whether real answer errors are well described by
\emph{any} per-attribute rate is the deepest of them, and only the study of
Appendix~\ref{app:protocol} can settle it. The channel is assumed
\emph{symmetric}, whereas real users may well have a bias---readier to say
``yes'' than ``no'' when unsure---which an asymmetric binary channel would
capture and which our formulation does not. And errors are correlated within an
attribute in Section~\ref{sec:results-correlated} but never \emph{across}
attributes, though a listener with an unusual notion of ``energetic'' plausibly
has an unusual notion of ``danceable'' too; that structure would penalise any
method that concentrates its questions within a tier, ours included.

\emph{One primary catalogue.} The headline comparison is on music. We add two
further catalogues and a synthetic family precisely because one catalogue
cannot support a general claim, and the resulting scope condition is stated as
a condition rather than as a universal.

\emph{Greedy selection.} We select greedily throughout. Greedy optimality under
noise is established for a single common crossover probability~\cite{jedynak2012}
and, as Remark~\ref{rem:greedy} notes, does not obviously extend to
heterogeneous per-attribute capacities. Holding the policy class fixed at the
one deployed systems use is what isolates the acquisition function; it is not a
claim that greedy is optimal here.

\emph{Static reliability.} $\hat\varepsilon$ is fixed for the session. A real
system could estimate a user's per-attribute reliability online from answer
consistency, which would make the specification problem largely disappear---and
would introduce an exploration/exploitation trade-off this paper does not
address.

\emph{Simulator--method coupling.} The simulator and the method share the
functional form of the noise model. This is standard in the noisy-search
literature but it is a real limitation: the ensemble experiment breaks the
parameter coupling but not the functional-form coupling, and the abstention
experiment is the only place where the true channel has a component the
baseline selector cannot represent.

\subsection{Future Work}

The obvious next step is the user study, and it is specified rather than
gestured at. Beyond it, three directions follow from what the experiments
showed rather than from what would be conventional to list.

\emph{Per-user} reliability estimation is the most valuable, and is the natural
extension of what Section~\ref{sec:results-em} already does. That estimator
pools every session into one crossover vector, which is the right object when
the population is homogeneous and the wrong one when it is not: a musician and
a casual listener do not share a crossover for ``what key is it in''. Estimating
$\hat\varepsilon_f$ per user, from within-session answer consistency and a
population prior, is a hierarchical version of the same EM and would attack the
variation this paper's model cannot express. It also introduces an
exploration/exploitation trade-off---questions asked to learn about the user
rather than about the item---that Section~\ref{sec:results-em} shows is not
free, and Section~\ref{sec:results-streaming} shows can be worse than not
free: we tried the simpler, non-hierarchical version of this idea---one
crossover vector, updated \emph{online} from the deployment's own trials
instead of offline from a fixed log---and it loses badly on the full
catalogue to a feedback loop the offline estimator's repeated replay does not
allow (\texttt{genre} monopolises $58.6\%$ of every question asked once the
running estimate correctly favours it, starving the other nineteen
attributes of the evidence needed to estimate them), though the deficit
disappears on a smaller catalogue with a longer budget. A per-user version
would face the identical loop at the scale of a single session rather than
the whole deployment, which is a reason for pessimism about a naive
implementation and a concrete requirement---replay, or an explicit
optimism bonus for under-sampled attributes, not exploration alone---for a
serious one.

Non-myopic selection under heterogeneous capacities is the most interesting
theoretically. Corollary~\ref{cor:budget} makes the difficulty concrete: with
per-attribute capacities, a horizon-optimal policy may wish to reserve a
high-capacity attribute for a later turn when the belief can exploit it, and a
greedy rule cannot express that preference.

Finally, the discretisation result suggests attacking the problem one stage
earlier. If quantile binning is what manufactures the bias, then a
reliability-aware \emph{discretiser}---one that chooses bin boundaries to
maximise mutual information through the channel rather than to equalise
populations---would prevent it rather than compensate for it. Nothing in this
paper requires the binning to be given.

%=============================================================================
\section{Conclusion}
\label{sec:conclusion}
%=============================================================================

Expected information gain is the default acquisition function for interactive
identification, and it was derived for an oracle that answers exactly. When the
oracle is a person, the criterion has a blind spot with a specific shape: it
scores a question by how evenly it divides the belief, and vagueness is what
makes a split even. In a catalogue that mixes quantile-binned perceptual
descriptors with skewed categorical metadata, the criterion therefore seeks out
the questions users answer worst---not as a failure of tuning but as a
consequence of what it optimises, sharpened by a preprocessing step that fixes
apparent informativeness at a value determined by the bin count rather than by
the attribute.

The correction is a single line, and it is textbook: score the question by the
mutual information between the target and the observed answer through a
per-attribute binary symmetric channel. What we have tried to establish is not
that this quantity is new but that it is needed, and that the knowledge it
requires is cheap enough to be had. Two coarse tiers recover most of the
benefit; the annotation may be wrong for a majority of attributes before it
stops paying; across sixty worlds drawn from priors on the tier values a system
shipping the fixed point estimates wins in sixty; and a two-line audit of the
catalogue predicts in advance whether any of this will help. The three numbers
need not even be assumed: EM over six hundred logged sessions of the
already-deployed incumbent, with the target item latent and nothing annotated,
recovers a crossover vector that matches the curator's exactly.

We have also been explicit about the boundaries. Reliability weighting is a
correction for a noisy channel, and on a clean one it is not merely useless but
harmful. It requires a reliability model, and inverting that model inverts the
result. The two routes that would have supplied one from the catalogue
alone---repeated-item instability, and a psychophysical model of the bin
boundaries---both fail, and fail for the same reason: they measure the data,
and answerability is a property of the people. The route that does work without
a curator is to watch those people, which is what the log-based estimator does
and why it succeeds where the catalogue-derived ones do not. And the parameter
has still not been measured on real people, only on simulated ones. That is the
piece of work this paper specifies and does not do, and it is the piece that
would close the argument.
